\documentclass[12pt]{article}

\input{../kz}

\rhead{Math 132H}

\DeclareMathOperator{\cis}{cis}
\DeclareMathOperator{\res}{res}
\newcommand{\lra}{\xLeftrightarrow}
\newcommand{\ra}{\xRightarrow}

\begin{document}

\section{Chapter 3}

\subsection{Exercise 21}

\subsubsection{Part A}

We'll prove the contrapositive.

Using a part of the proof from the five equivalences in the notes,
we know that $\Omega$ not being simply connected implies that
we can separate it into $\Omega = K \cup L$, where $\infty \in L$ and
$K$ and $L$ are both open in $K \cup L$.

Take any $k \in K$ and consider the vertical line $x=\Re(k)$.

Notice that this line must cross $\Omega$ both below and above $k$.
If any of them didn't intersect, then that line would go off to $\infty$
and connect $k$ with $L$, violating the separation of the two.

Thus, $\exists a, b \in \Omega: \Re(a)=\Re(b)=\Re(k) \land \Im(a) < \Im(k) < \Im(k)$.

By high school algebra, $k$ lies on the line segment directly connecting $a$ and $b$,
banning $\Omega$ from being convex.

\subsubsection{Part B}

Once again, take our $K$, $L$, and $k$ like in the previous part.

Now consider any point $z_0 \in \Omega$ and the line from $z_0$ to $k$.
If we extend this line, it must intersect $\Omega$ at some other point $z$,
as otherwise $K$ would be connected to $L$.

Since $k$ lies on the line segment from $z_0$ to $z$, $z_0$ cannot serve as the
center for our star-shaped domain.
The choice of $z_0$ was arbitrary, so $\Omega$ in its entirety is not star-shaped.

\subsection{Exercise 22}

Consider the function $g(z)=z \cdot f(z)$.

We can write $g(0)=0$ in the form of a Cauchy integral
(at least I \textit{think} that's what this is):
\begin{align*}
  g(0)
  &=\frac{1}{2\pi i} \oint_{|z|=r} \frac{g(z)}{z}\,dz \\
  &= \frac{1}{2\pi i} \oint_{|z|=r} f(z)\,dz
\end{align*}
The integral is continuous for $r \le 1$, but when $r=1$ $\oint_{\partial \mathbb{D}} f(z)\,dz=2\pi i$,
which is nonzero and contradicts our earlier statement that $g(0)=0$.

\pagebreak

\section{Chapter 8}

\subsection{Exercise 1}

\subsubsection{Forwards}

I'll show the contrapositive: assume a $z_0 \in U: f'(z_0)=0$.
WLOG assume $f(z_0)=0$ as well.

We can shrink a disc around $z_0$ to a point where $f$ and $f'$ only have zeros at $z_0$.
Let the a radius at which this happens be denoted as $R$.

I'll show that for any $0 < r < R$, $f$ is not injective in $\{z: |z-z_0| < r\}$.
For convenience, I'll let $B$ define the border of this disc.

We have by our premise that $f$ has two zeros in the disc.
Also, $\inf_{z \in B} f(z)=\epsilon > 0$, since circles are compact.
Thus, if we consider any $w: 0 < |w| < \epsilon$, $|f(z)| > |w|$ on $B$
and we can apply Rouche's Theorem to see that $f(z)-w$ has two zeros.

Notice that these two zeros (call them $z_1$ and $z_2$) can't be dupes.
If it were, then $\frac{d}{dz} (f(z_1)-w)=0$ and we know $f'$ is only zero at $z_0$.

Thus, no matter how small we shrink our radius, we can always find two distinct
points that get sent by $f$ to the same destination.

\subsubsection{Backwards}

Fix a $z_0 \in U$ and WLOG let $f(z_0)=0$.

Consider a disc of radius $r$ around $z_0$ small enough that $z_0$ is the only zero in it.
We can do this since otherwise it would imply a limit point and $f$ being constant.

In this disc, we can write $f(z)=f'(z_0)(z-z_0)+o(|z-z_0|)$.

Now, I lied, since the disc we're actually going to use the disc of radius $\frac{r}{2}$ as our candidate.

Fix any $x$ in this smaller disc.
On the border of the \textit{larger} disc,
\begin{align*}
  |f'(z_0)(z-z_0)-f(x)|
  &= |f'(z_0)(z-z_0)-f'(z_0)(x-z_0)+o(|x-z_0|)| \\
  &= |f'(z_0)(z-x)+o(|x-z_0|)| \\
  &> \frac{r}{2}f'(z_0) - |o(|x-z_0|)| \\
  &> |o(|z-z_0|)|
\end{align*}
where the last bound becomes true if we retroactively shrink $r$ enough.

But yeah, $f'(z_0)(z-z_0)-f(x)$ has only one zero in either disc.

Using Rouche's Theorem, we get that $(f'(z_0)(z-z_0)-f(x))+(o(|z-z_0|))=f(z)-f(x)$ must have only one as well.

\pagebreak

\subsection{Exercise 3}

I'll use the textbook definition of simply connected, which means
that every curve can be continuously deformed to another curve.

We have a continous bijection $f$ between $U$ and $V$.

Now consider any two curves $\gamma_1$ and $\gamma_2$ in $V$.
Through $f$, we can obtain their preimages in $U$, $\varphi_1$ and $\varphi_2$.
Since $f$ is continuous, these are also curves.

$U$ is simply connected, so $\exists d: [0, 1] \times [0, 1] \to U$
where $\varphi_1(t)=d(0, t)$ and $\varphi_2(t)=d(1, t)$.

Since $f$ and $d$ are both continous, their composition $f \circ d$ is too.
Notice that this is exactly the deformation we want, since the first argument
being $0$ corresponds to $\gamma_1$ and it being $1$ corresponds to $\gamma_2$.

Thus, $V$ must also be simply connected. $\square$

\subsection{Exercise 9}

The function $f(z)=\frac{i+z}{i-z}$ is holomorphic everywhere except $i$.

$u$ is the real part of $f$, so it must be harmonic where $f$ is holomorphic,
and that includes $\mathbb{D}$.

To show that $u$ vanishes everywhere on $\partial \mathbb{D}$ is just calculations:
\begin{align*}
  \frac{i+e^{i\theta}}{i-e^{i\theta}}
  &= \frac{i+\cos\theta+i\sin\theta}{i-\cos\theta-i\sin\theta} \\
  &= \frac{(i+\cos\theta+i\sin\theta)(-\cos\theta-i+i\sin\theta)}{(i-\cos\theta-i\sin\theta)(-\cos\theta-i+i\sin\theta)} \\
  &= \frac{1-\cos^2\theta-\sin^2\theta+i(\cos\theta(1-\sin\theta)-\cos\theta(1+\sin\theta))}{c} \\
  &= \frac{-2i\sin\theta\cos\theta}{c}
\end{align*}
I've turned the denominator into some $c \in \R$ since it doesn't really matter what it is.

These equations break when $\theta=\frac{\pi}{2}$, but it's fine
since the problem manually defines $u(i)=0$.

Regardless, $u(z)$ is indeed $0$ on the border. $\square$

\end{document}
